\documentclass[11pt,letterpaper]{article}

\usepackage[utf8]{inputenc}
\usepackage[T1]{fontenc}
\usepackage[spanish,es-nodecimaldot]{babel}
\usepackage{lmodern}
\usepackage{geometry}
\usepackage{booktabs}
\usepackage{longtable}
\usepackage{tabularx}
\usepackage{array}
\usepackage{float}
\usepackage{amsmath}
\usepackage{xcolor}
\usepackage{listings}
\usepackage{fancyhdr}
\usepackage{hyperref}
\usepackage{pdflscape}
\usepackage{enumitem}
\usepackage{setspace}

\geometry{margin=2.35cm}
\setlength{\parindent}{0pt}
\setlength{\parskip}{0.45em}
\setstretch{1.10}
\setcounter{tocdepth}{2}
\setcounter{secnumdepth}{3}

% ============================================================
% DATOS DE PORTADA
% ============================================================
\newcommand{\Universidad}{Instituto Tecnológico de Costa Rica}
\newcommand{\Escuela}{Escuela de Ingeniería en Computadores}
\newcommand{\Curso}{CE 1108 -- Compiladores e Intérpretes}
\newcommand{\TituloProyecto}{Proyecto: Compilador para Código Binario}
\newcommand{\Profesor}{Jason Leitón Jiménez}
\newcommand{\FechaEntrega}{14 de octubre de 2026, 13:00}
\newcommand{\Repositorio}{https://github.com/FabiMC040105/CompiladoresIIS2026.git}

% Nombres y carnés tomados del Avance 1
\newcommand{\EstudianteUno}{Camila Lizano Brenes}
\newcommand{\CarneUno}{2024255324}
\newcommand{\EstudianteDos}{Fabiana Moreno Castrillo}
\newcommand{\CarneDos}{2023153700}
\newcommand{\EstudianteTres}{Felipe Viales Vieto}
\newcommand{\CarneTres}{2021030084}
\newcommand{\EstudianteCuatro}{Daniel González Chaves}
\newcommand{\CarneCuatro}{2022279667}

% ============================================================
% ESTILOS
% ============================================================
\definecolor{codigoFondo}{RGB}{248,248,248}
\definecolor{codigoBorde}{RGB}{210,210,210}
\definecolor{palabra}{RGB}{35,65,150}
\definecolor{comentario}{RGB}{70,115,70}

\lstdefinestyle{codigo}{
    basicstyle=\ttfamily\footnotesize,
    backgroundcolor=\color{codigoFondo},
    frame=single,
    rulecolor=\color{codigoBorde},
    breaklines=true,
    columns=fullflexible,
    keepspaces=true,
    showstringspaces=false,
    keywordstyle=\bfseries\color{palabra},
    commentstyle=\color{comentario}
}
\lstdefinestyle{asm}{style=codigo,numbers=none}
\lstdefinestyle{terminal}{
    style=codigo,
    numbers=none,
    basicstyle=\ttfamily\scriptsize
}

% Inserta evidencia externa si existe. Si todavía no existe, deja un marcador visible.
\newcommand{\ArchivoOpcional}[2]{%
    \IfFileExists{#1}{%
        \lstinputlisting[style=#2]{#1}%
    }{%
        \begin{center}
        \fbox{\parbox{0.92\linewidth}{\textbf{Pendiente de evidencia final:} \texttt{\detokenize{#1}}.\\
        Copiar aquí el archivo generado por la versión final antes de entregar.}}
        \end{center}
    }%
}

\hypersetup{
    colorlinks=true,
    linkcolor=blue!50!black,
    urlcolor=blue!60!black,
    pdftitle={Especificación de la Arquitectura y Documentación del Compilador},
    pdfauthor={Grupo CE1108}
}

\pagestyle{fancy}
\fancyhf{}
\lhead{CE 1108 -- Compiladores e Intérpretes}
\rhead{Compilador para Código Binario}
\cfoot{\thepage}
\setlength{\headheight}{14pt}

\begin{document}

% ============================================================
% PORTADA
% ============================================================
\begin{titlepage}
    \thispagestyle{empty}
    \centering
    \vspace*{1.0cm}

    {\Large\bfseries \Universidad\par}
    \vspace{0.15cm}
    {\large \Escuela\par}

    \vspace{1.4cm}

    {\Large\bfseries \Curso\par}
    \vspace{0.7cm}
    {\LARGE\bfseries \TituloProyecto\par}
    \vspace{0.35cm}
    {\large Especificación de la Arquitectura y Documentación del Compilador\par}

    \vspace{1.25cm}

    \renewcommand{\arraystretch}{1.25}
    \begin{tabular}{>{\bfseries}p{8.0cm}p{3.0cm}}
        Estudiante & Carné \\
        \midrule
        \EstudianteUno & \CarneUno \\
        \EstudianteDos & \CarneDos \\
        \EstudianteTres & \CarneTres \\
        \EstudianteCuatro & \CarneCuatro \\
    \end{tabular}

    \vfill

    \begin{tabular}{>{\bfseries}ll}
        Profesor: & \Profesor \\
        Fecha de entrega: & \FechaEntrega \\
    \end{tabular}

    \vspace{0.45cm}

    {\small\textbf{Repositorio del proyecto:} \url{\Repositorio}\par}

    \vspace{0.8cm}
\end{titlepage}

\tableofcontents
\clearpage

% ============================================================
% SECCIÓN 3 DEL ENUNCIADO
% ============================================================
\setcounter{section}{2}
\section{Especificación de la Arquitectura}
\label{sec:arquitectura}

La arquitectura objetivo corresponde al procesador VLIW desarrollado por el grupo en Arquitectura de Computadores I. El compilador genera ensamblador para esta ISA y debe respetar tanto su codificación como las restricciones de ejecución del pipeline.

\subsection{Tamaño de palabra}

Los datos se manejan en palabras de 32 bits. Cada instrucción VLIW completa, denominada \textit{bundle}, tiene 128 bits y contiene cuatro slots fijos de 32 bits:

\begin{center}
\begin{tabular}{|c|c|c|c|}
\hline
\textbf{FEIST [127:96]} & \textbf{BRU [95:64]} & \textbf{LSU [63:32]} & \textbf{ALU [31:0]} \\
\hline
32 bits & 32 bits & 32 bits & 32 bits \\
\hline
\end{tabular}
\end{center}

El slot ALU ocupa los bits menos significativos y FEIST los más significativos. Los campos no utilizados se rellenan con cero. Un bundle ocupa 16 bytes, por lo que el PC avanza 16 bytes en ejecución secuencial. El orden de bytes es \textit{little-endian}.

\subsection{Conjunto de registros}

La arquitectura dispone de 16 registros identificados de \texttt{x0} a \texttt{x15}. El número del registro es también su código de 4 bits dentro de las instrucciones.

\begin{table}[H]
\centering
\caption{Registros y convención utilizada por el compilador.}
\begin{tabularx}{\textwidth}{>{\ttfamily}p{1.8cm}X}
\toprule
\textnormal{Registro} & \textnormal{Propósito} \\
\midrule
x0 & Constante cero; las escrituras se descartan. \\
x1 & Program Counter (PC). Leerlo entrega el PC del bundle; una escritura válida redirige el flujo. \\
x2 & Valor de retorno de funciones por convenio del backend. \\
x3 & Registro auxiliar reservado por el backend. \\
x4--x12 & Registros temporales. \\
x13 & Registro auxiliar para direcciones y cálculos de memoria. \\
x14 & Stack Pointer (SP). \\
x15 & Frame Pointer (FP). \\
\bottomrule
\end{tabularx}
\end{table}

\subsection{Organización de memoria}

La memoria de datos implementada es de 64 KiB. El compilador la divide lógicamente de la siguiente forma:

\begin{table}[H]
\centering
\caption{Organización lógica de la memoria de datos.}
\begin{tabular}{lll}
\toprule
\textbf{Región} & \textbf{Rango} & \textbf{Uso} \\
\midrule
Datos globales & \texttt{0x0000--0x7FFF} & Variables globales \\
Pila & \texttt{0x8000--0xFFFF} & Parámetros, marcos y variables locales \\
\bottomrule
\end{tabular}
\end{table}

Los escalares \texttt{entero} y \texttt{vof} ocupan 4 bytes. Los elementos de listas y matrices también ocupan 4 bytes. Las variables globales se asignan desde direcciones bajas hacia arriba, mientras que la pila crece hacia direcciones menores. El SP se inicializa conceptualmente en \texttt{0x10000}; la primera reserva lo desplaza a una dirección válida dentro de la región de pila. Los accesos \texttt{LOAD} y \texttt{STORE} son de 32 bits y se realizan sobre direcciones alineadas a 4 bytes.

La memoria de instrucciones está separada de la memoria de datos. Cada bundle consume 16 bytes; con una memoria de instrucciones de 64 KiB pueden almacenarse hasta 4096 bundles. La versión actual no utiliza \textit{heap} ni asignación dinámica.

\subsection{Conjunto de instrucciones}

\begin{landscape}
\begin{longtable}{p{1.7cm}p{1.5cm}p{1.5cm}p{5.4cm}p{1.5cm}p{7.2cm}}
\caption{Tabla completa de instrucciones de la ISA.}\label{tab:isa}\\
\toprule
\textbf{Instr.} & \textbf{Slot} & \textbf{Formato} & \textbf{Sintaxis} & \textbf{Opcode} & \textbf{Descripción} \\
\midrule
\endfirsthead
\toprule
\textbf{Instr.} & \textbf{Slot} & \textbf{Formato} & \textbf{Sintaxis} & \textbf{Opcode} & \textbf{Descripción} \\
\midrule
\endhead
NOP & Todos & -- & \texttt{NOP} & 0000 & No realiza operación; el slot se representa con ceros. \\
ADD & ALU & R & \texttt{ADD rd, rs1, rs2} & 0001 & Suma de 32 bits. \\
SUB & ALU & R & \texttt{SUB rd, rs1, rs2} & 0010 & Resta de 32 bits. \\
OR & ALU & R & \texttt{OR rd, rs1, rs2} & 0011 & OR bit a bit. \\
XOR & ALU & R & \texttt{XOR rd, rs1, rs2} & 0100 & XOR bit a bit. \\
AND & ALU & R & \texttt{AND rd, rs1, rs2} & 0101 & AND bit a bit. \\
SLT & ALU & R & \texttt{SLT rd, rs1, rs2} & 0110 & Escribe 1 si \texttt{rs1 < rs2} con signo. \\
ADDI & ALU & I & \texttt{ADDI rd, rs1, imm} & 1001 & Suma un inmediato con signo de 12 bits. \\
SRLI & ALU & I & \texttt{SRLI rd, rs1, imm} & 1010 & Desplazamiento lógico a la derecha. \\
SLLI & ALU & I & \texttt{SLLI rd, rs1, imm} & 1011 & Desplazamiento lógico a la izquierda. \\
LOAD & LSU & L & \texttt{LOAD rd, off(rb)} & 01 & Carga una palabra de memoria. \\
STORE & LSU & S & \texttt{STORE rs, off(rb)} & 10 & Almacena una palabra en memoria. \\
JMP & BRU & J & \texttt{JMP offset} & 001 & Salto incondicional PC-relativo. \\
BGE & BRU & B & \texttt{BGE rs1, rs2, offset} & 010 & Salta si \texttt{rs1 >= rs2}. \\
BLE & BRU & B & \texttt{BLE rs1, rs2, offset} & 011 & Salta si \texttt{rs1 <= rs2}. \\
BEQ & BRU & B & \texttt{BEQ rs1, rs2, offset} & 100 & Salta si los operandos son iguales. \\
BNE & BRU & B & \texttt{BNE rs1, rs2, offset} & 101 & Salta si los operandos son diferentes. \\
AUTH & FEIST & FA & \texttt{AUTH rs1, rs2} & 001 & Intenta autorizar el acceso a la bóveda. \\
LOGOUT & FEIST & FL & \texttt{LOGOUT} & 010 & Elimina la autorización actual. \\
KSTORE & FEIST & FK & \texttt{KSTORE key, round, rs1} & 011 & Guarda una subllave en la bóveda. \\
FROUND & FEIST & FR & \texttt{FROUND rdL, rdR, rsL, rsR, key, round, mode} & 100 & Ejecuta una ronda de Feistel4. \\
\bottomrule
\end{longtable}
\end{landscape}

Ejemplos de uso:

\begin{lstlisting}[style=asm]
ALU: ADD x4, x2, x3
ALU: ADDI x5, x4, 12
LSU: LOAD x6, -4(x15)
LSU: STORE x6, -8(x15)
BRU: BNE x4, x0, 64
BRU: JMP -160
FEIST: AUTH x2, x3
FEIST: KSTORE 0, 2, x4
FEIST: FROUND x7, x8, x5, x6, 0, 2, 0
FEIST: LOGOUT
\end{lstlisting}

Las operaciones de alto nivel \texttt{MUL}, \texttt{DIV}, módulo, potencia, llamadas y retornos no tienen un opcode nativo propio; el compilador las sintetiza mediante secuencias de instrucciones disponibles.

\subsubsection{Formatos principales}

\textbf{ALU tipo R:}
\[
[31{:}16]=0,\quad [15{:}12]=rs2,\quad [11{:}8]=rs1,\quad [7{:}4]=rd,\quad [3]=0,\quad [2{:}0]=opcode
\]

\textbf{ALU tipo I:}
\[
[31{:}24]=0,\quad [23{:}12]=imm,\quad [11{:}8]=rs1,\quad [7{:}4]=rd,\quad [3]=1,\quad [2{:}0]=opcode
\]

\textbf{LSU:}
\[
[31{:}22]=0,\quad [21{:}10]=offset,\quad [9{:}6]=rb,\quad [5{:}2]=reg,\quad [1{:}0]=opcode
\]

\textbf{BRU condicional:}
\[
[31{:}23]=0,\quad [22{:}11]=offset,\quad [10{:}7]=rs2,\quad [6{:}3]=rs1,\quad [2{:}0]=opcode
\]

\textbf{JMP:}
\[
[31{:}23]=0,\quad [22{:}3]=offset,\quad [2{:}0]=001
\]

Los offsets de BRU se expresan en bytes y se suman al PC del bundle que contiene el salto. Los condicionales utilizan 12 bits con signo y \texttt{JMP} utiliza 20 bits con signo.

% ============================================================
% SECCIÓN 8 DEL ENUNCIADO
% ============================================================
\clearpage
\setcounter{section}{7}
\section{Documentación del Compilador}

El código fuente, las pruebas y los artefactos generados del proyecto se encuentran en el repositorio oficial del grupo: \url{\Repositorio}.

\subsection{Especificación de la Arquitectura}

La especificación solicitada en este apartado corresponde a la misma arquitectura objetivo documentada en la Sección~\ref{sec:arquitectura}. Allí se incluyen el tamaño de palabra e instrucción, registros y propósitos, organización de memoria, tabla completa del ISA, opcodes, formatos y ejemplos de uso.

\subsection{Especificación del Lenguaje de Entrada}

El lenguaje fuente utiliza la extensión \texttt{.e}. Las sentencias terminan con la secuencia \texttt{:)} y los bloques se delimitan con llaves.

\subsubsection{Tokens}

\begin{table}[H]
\centering
\caption{Resumen de tokens del lenguaje.}
\begin{tabularx}{\textwidth}{>{\bfseries}p{3.1cm}X}
\toprule
Categoría & Elementos reconocidos \\
\midrule
Palabras reservadas & \texttt{si}, \texttt{sinosi}, \texttt{sino}, \texttt{mientras}, \texttt{para}, \texttt{defina}, \texttt{devuelva}, \texttt{rompa}, \texttt{principal}. \\
Tipos & \texttt{entero}, \texttt{vof}, \texttt{vacio}. \\
Colecciones & \texttt{lista}, \texttt{matriz}. \\
Booleanos & \texttt{verdadero}, \texttt{falso}. \\
Aritméticos & \texttt{+}, \texttt{-}, \texttt{*}, \texttt{/}, \texttt{\%}, \texttt{**}. \\
Comparación & \texttt{==}, \texttt{!=}, \texttt{<}, \texttt{>}, \texttt{<=}, \texttt{>=}. \\
Lógicos & \texttt{\$\$} (AND) y \texttt{--} (OR). \\
Delimitadores & \texttt{\{ \} ( ) [ ] ,} y terminador \texttt{:)}. \\
Identificadores & Deben iniciar con letra minúscula; luego pueden contener letras, dígitos o \texttt{\_}. \\
Números & \texttt{0} o enteros positivos sin ceros iniciales; el signo negativo se maneja como operador. \\
\bottomrule
\end{tabularx}
\end{table}

Los comentarios de una línea inician con el carácter 	exttt{\char94}. Los comentarios de bloque se delimitan con dos caracteres 	exttt{\char94\char94}.

\begin{lstlisting}[style=codigo]
^ comentario de una linea

^^
comentario de varias lineas
^^
\end{lstlisting}

\subsubsection{Variables, asignaciones y expresiones}

\begin{lstlisting}[style=codigo]
entero contador = 0 :)
vof activo = verdadero :)
contador = contador + 1 :)
activo = contador < 10 :)
\end{lstlisting}

Las listas y matrices son estáticas. Cada elemento ocupa 4 bytes y las matrices se almacenan por filas (\textit{row-major}):

\begin{lstlisting}[style=codigo]
entero lista numeros[4] = [10, 20, 30, 40] :)
entero matriz tabla[2][2] = [[1, 2], [3, 4]] :)
tabla[1][0] = numeros[2] :)
\end{lstlisting}

Para una matriz con $C$ columnas:
\[
\operatorname{dir}(i,j)=\operatorname{base}+4(iC+j)
\]

\subsubsection{Estructuras de control y funciones}

\begin{lstlisting}[style=codigo]
si (x > 10) {
    x = x - 1 :)
} sinosi (x == 10) {
    x = x + 2 :)
} sino {
    x = x + 1 :)
} :)

mientras (i < 10) {
    i = i + 1 :)
} :)

para (entero j = 0, j < 5, j = j + 1) {
    x = x + j :)
} :)

defina entero sumar(entero a, entero b) {
    devuelva a + b :)
} :)
\end{lstlisting}

La sentencia \texttt{rompa :)} termina el ciclo más interno. Las funciones pueden recibir parámetros y devolver valores mediante \texttt{devuelva}. El compilador controla ámbitos globales, de función y de bloque; por ello, dos identificadores con el mismo nombre pueden representar símbolos distintos si pertenecen a ámbitos diferentes.

\subsection{Descripción de Algoritmos Clave}

\subsubsection{Análisis léxico}

El archivo \texttt{lexer.l} implementa el analizador con Flex. Cada regla asocia un patrón con un token de Bison. Los espacios y comentarios se ignoran; los lexemas válidos retornan su token y los caracteres no reconocidos producen error léxico. Para comentarios de bloque se utiliza un estado exclusivo, lo cual evita que su contenido sea interpretado como código.

\subsubsection{Análisis sintáctico y AST}

El archivo \texttt{parser.y} define la gramática en Bison. Cuando una producción es reconocida se construyen nodos del árbol de sintaxis abstracta (AST). Cada nodo conserva su tipo, valor cuando corresponde, línea de origen, primer hijo y siguiente hermano. Esto permite representar bloques y secuencias de instrucciones sin almacenar detalles sintácticos innecesarios.

\subsubsection{Tabla de símbolos y memoria}

La tabla de símbolos registra nombre, categoría, tipo, forma, ámbito y ubicación de memoria. La búsqueda de identificadores comienza en el ámbito actual y continúa por sus ámbitos padre. Las variables globales reciben direcciones absolutas; los parámetros y locales utilizan desplazamientos respecto a \texttt{x15} (FP).

La convención principal del marco de función es:

\begin{center}
\begin{tabular}{ll}
\toprule
\textbf{Posición} & \textbf{Contenido} \\
\midrule
\texttt{FP+8, FP+12, ...} & Parámetros \\
\texttt{FP+4} & Dirección de retorno \\
\texttt{FP+0} & FP anterior \\
\texttt{FP-4, FP-8, ...} & Variables locales \\
\bottomrule
\end{tabular}
\end{center}

\subsubsection{Generación de ensamblador y riesgos de datos}

El backend recorre el AST y emite instrucciones ALU, LSU, BRU y FEIST. Como el procesador usa el pipeline

\begin{center}
\texttt{IF $\rightarrow$ ID $\rightarrow$ RR $\rightarrow$ EX $\rightarrow$ WB},
\end{center}

el compilador calendariza las dependencias RAW. Si un productor se emite en el bundle $B$, el valor puede consumirse de forma segura a partir de $B+3$. Por ello, una dependencia inmediata requiere 2 NOPs; si ya existe un bundle independiente requiere 1; y con dos bundles independientes no requiere NOP adicional.

\begin{lstlisting}[style=asm]
ALU: ADDI x4, x0, 5
ALU: NOP
ALU: NOP
ALU: ADD x5, x4, x6
\end{lstlisting}

Las escrituras a \texttt{x1} se tratan de forma especial porque redirigen el PC. El backend agrega una barrera de control de un bundle después del retorno por software.

\subsubsection{Cálculo de saltos y resolución de etiquetas}

Las estructuras de control se generan inicialmente con etiquetas simbólicas. El resolvedor realiza dos pasos:

\begin{enumerate}[leftmargin=*]
    \item recorre el ensamblador y registra cada etiqueta junto con el índice de bundle en el que se encuentra;
    \item recorre las referencias, localiza su destino y sustituye el nombre por un desplazamiento numérico.
\end{enumerate}

Como cada bundle ocupa 16 bytes, para una referencia desde el bundle $B_o$ hacia $B_d$ se calcula:

\[
\text{offset}=(B_d-B_o)\times 16
\]

La resolución se hace después de insertar los NOPs, ya que estos cambian la posición definitiva de los saltos. El módulo también detecta etiquetas duplicadas, referencias inexistentes y desplazamientos fuera del rango representable por la instrucción.

\subsubsection{Resolución de referencias a funciones}

Las llamadas y retornos se implementan por software. El llamador preserva temporales, coloca argumentos en la pila y salta a la etiqueta de la función. El valor retornado se recibe en \texttt{x2}. La dirección de retorno se maneja simbólicamente durante la generación y se convierte posteriormente a un desplazamiento respecto al PC. El callee restaura el marco anterior y escribe la dirección de retorno en \texttt{x1} para redirigir la ejecución al llamador.

\subsection{Ejemplos de Ejecución}
\label{sec:ejemplos_ejecucion}

Para validar el flujo completo del compilador se utilizaron los tres programas requeridos por el proyecto. En cada caso se conserva el código fuente, el ensamblador generado, la representación hexadecimal de los bundles VLIW y una traza de ejecución obtenida sobre el procesador de Arquitectura de Computadores I. Los archivos de evidencia se almacenan en \texttt{docs/evidencias/}.

La representación hexadecimal contiene una línea de 32 dígitos hexadecimales por bundle, equivalentes a 128 bits. Para la validación funcional se cargó cada archivo en la memoria de instrucciones del procesador y se ejecutó con el banco de pruebas de integración \texttt{tb\_compilador.sv}. La llegada al bucle terminal \texttt{fin\_programa}, junto con los valores finales de registros y memoria, se utilizó como criterio de éxito.

\begin{table}[H]
\centering
\caption{Resumen de resultados de los tres programas de validación.}
\begin{tabularx}{\textwidth}{>{\bfseries}p{3.2cm}X p{2.4cm} p{2.2cm}}
\toprule
Programa & Resultado validado & PC de \texttt{fin\_programa} & Estado \\
\midrule
Factorial & $5! = 120$; el valor final se observa en \texttt{x2} y en memoria. & \texttt{0x00000F70} & PASS \\
Búsqueda en matriz & Para el objetivo 23 se obtiene \texttt{posicion=4}, \texttt{fila=1} y \texttt{columna=1}. & \texttt{0x00001950} & PASS \\
MCD & $\operatorname{MCD}(48,18)=6$; el valor retornado queda en \texttt{x2}. & \texttt{0x00002000} & PASS \\
\bottomrule
\end{tabularx}
\end{table}

\subsubsection{Ejercicio 1: cálculo de factorial}

El primer programa calcula el factorial de 5 mediante una función, variables locales, un condicional y un ciclo. La ejecución completa produjo el valor 120 y alcanzó correctamente \texttt{fin\_programa} en \texttt{0x00000F70}.

\paragraph{Código fuente completo.}
\ArchivoOpcional{evidencias/factorial.e}{codigo}

\paragraph{Salida de compilación.}
\ArchivoOpcional{evidencias/factorial_compilacion.txt}{terminal}

\paragraph{Ensamblador intermedio generado.}
\ArchivoOpcional{evidencias/factorial.asm}{asm}

\paragraph{Representación hexadecimal legible.}
\ArchivoOpcional{evidencias/factorial.hex}{asm}

\paragraph{Traza de ejecución.}
En la traza se observa la evolución del resultado parcial hasta 120, la restauración del marco de pila y la detección del bucle terminal. En particular, \texttt{x2=120} corresponde al valor retornado por la función factorial.

\ArchivoOpcional{evidencias/traza_factorial.txt}{terminal}

\clearpage
\subsubsection{Ejercicio 2: búsqueda en matriz bidimensional}

El segundo programa declara la matriz

\begin{center}
$\left[\begin{array}{ccc}
4 & 8 & 15 \\
16 & 23 & 42
\end{array}\right]$
\end{center}

y busca el valor 23 mediante dos ciclos anidados. La función devuelve la posición lineal en orden por filas. El resultado validado fue \texttt{posicion=4}; posteriormente el programa calcula \texttt{fila=1} y \texttt{columna=1}. El procesador alcanzó \texttt{fin\_programa} en \texttt{0x00001950}.

\paragraph{Código fuente completo.}
\ArchivoOpcional{evidencias/ejercicio2_matriz.e}{codigo}

\paragraph{Salida de compilación.}
\ArchivoOpcional{evidencias/ejercicio2_matriz_compilacion.txt}{terminal}

\paragraph{Ensamblador intermedio generado.}
\ArchivoOpcional{evidencias/ejercicio2_matriz.asm}{asm}

\paragraph{Representación hexadecimal legible.}
\ArchivoOpcional{evidencias/ejercicio2_matriz.hex}{asm}

\paragraph{Traza de ejecución.}
La traza final verifica los valores de las variables principales en memoria: \texttt{objetivo=23}, \texttt{posicion=4}, \texttt{fila=1} y \texttt{columna=1}. También muestra los saltos tomados durante los ciclos y el bucle terminal del programa.

\ArchivoOpcional{evidencias/traza_matriz.txt}{terminal}

\clearpage
\subsubsection{Ejercicio 3: máximo común divisor}

Como programa propuesto se utilizó el cálculo del máximo común divisor de 48 y 18. El programa emplea funciones con parámetros y retorno, incluyendo una función auxiliar de valor absoluto, llamadas entre funciones y estructuras de control para implementar el algoritmo. La ejecución produjo \texttt{MCD(48,18)=6} y alcanzó \texttt{fin\_programa} en \texttt{0x00002000}.

\paragraph{Código fuente completo.}
\ArchivoOpcional{evidencias/ejercicio3_mcd.e}{codigo}

\paragraph{Salida de compilación.}
\ArchivoOpcional{evidencias/ejercicio3_mcd_compilacion.txt}{terminal}

\paragraph{Ensamblador intermedio generado.}
\ArchivoOpcional{evidencias/ejercicio3_mcd.asm}{asm}

\paragraph{Representación hexadecimal legible.}
\ArchivoOpcional{evidencias/ejercicio3_mcd.hex}{asm}

\paragraph{Traza de ejecución.}
Durante la ejecución se observan los valores intermedios del algoritmo hasta llegar a 6. El valor final queda en \texttt{x2=6}, y la traza termina con la detección del bucle de \texttt{fin\_programa}.

\ArchivoOpcional{evidencias/traza_mcd.txt}{terminal}

\paragraph{Resultado de la integración.}
Los tres programas fueron compilados por la versión final del compilador y ejecutados en el procesador VLIW de CE4301. En los tres casos se alcanzó el PC terminal esperado y los valores observados coincidieron con el resultado semántico esperado. Esta prueba valida de extremo a extremo el flujo \texttt{fuente .e $\rightarrow$ ensamblador $\rightarrow$ HEX VLIW $\rightarrow$ ejecución RTL}.

\subsection{Manejo de Casos Especiales}

\subsubsection{Saltos anidados}

Cada estructura de control genera etiquetas únicas. Un contador interno evita colisiones entre etiquetas pertenecientes a distintos \texttt{si}, \texttt{mientras} o \texttt{para}. La sentencia \texttt{rompa} utiliza la etiqueta de salida del ciclo más interno activo. Al resolver las referencias se calcula la posición definitiva de cada etiqueta después de la calendarización.

\subsubsection{Llamadas a función y retorno de valores}

El compilador utiliza \texttt{x14} como SP, \texttt{x15} como FP y \texttt{x2} para el valor retornado. Los argumentos se colocan en la pila y el callee crea un nuevo marco. El retorno no depende de instrucciones \texttt{CALL}/\texttt{RET}, ya que la ISA no las posee; se implementa mediante \texttt{JMP} para la llamada y una escritura a \texttt{x1} para regresar.

\subsubsection{Variables con el mismo nombre en distintos ámbitos}

Cada símbolo conserva el identificador del ámbito donde fue declarado. La resolución busca primero en el ámbito actual y luego asciende por la jerarquía. Por ello, una variable local puede ocultar temporalmente una variable externa con el mismo nombre sin alterar la entrada correspondiente al ámbito superior.

\subsubsection{Asignación de memoria local y global}

Las variables globales utilizan direcciones absolutas dentro de \texttt{0x0000--0x7FFF}. Los parámetros y variables locales se ubican en el marco de la función y se acceden mediante desplazamientos respecto a FP. Las listas y matrices reservan espacio contiguo para todos sus elementos; cuando se pasan como parámetros se utiliza una dirección de 4 bytes.

\vfill
\begin{center}
\small
\textit{Las evidencias de la Sección 8.4 corresponden a la versión final validada del compilador y a ejecuciones reales sobre el procesador VLIW.}
\end{center}

\end{document}
